%% ch11 --- Fractals: measuring what has no length
%%
%% Written September 2026. Every number the reader cannot do in their head
%% comes from code/measure/ch11.py, which imports the SAME functions the
%% listings print (code/ch11/), so the page and the listing cannot disagree.
%% The Henon counts use only +, - and * and boxes whose sides are powers of
%% two, so they are identical on every machine; dimensions are committed to
%% two decimals.
%%
%% WHAT THIS CHAPTER MAY NOT SAY. It measures the Henon attractor and says
%% that strange attractors are fractals; that the attractor has zero area and
%% infinite length is Chapter 10's, and is quoted rather than re-derived. The
%% Mandelbrot set and Julia sets are Chapter 12's: this chapter names the
%% set in one closing sentence and draws nothing in the complex plane.
%%
%% Twin: chapters/pl/ch11-fractals.tex. Section for section, macro for macro,
%% number for number; tools/parity.py compares them.

\chapter{Fractals: measuring what has no length}
\label{ch:fractals}
\index{fractal}\index{coastline paradox}

\noindent
How long is the coast of Britain? It sounds like a question for an atlas.
Look it up in two reference books, though, and you may well find two figures
too far apart for either to be a rounding of the other, and neither of them
a misprint. Why?

This chapter answers that, and then replaces length with a number that does
describe shapes like a coast. You will build two such shapes by repeating
one rule, measure the new number by counting boxes, and then measure the
shape chaos left behind in Chapter~\ref{ch:strange} the same way.

\begin{outcomes}
\outcome{measure a curve with rulers of different lengths and say whether its
  length settles}
\outcome{build the Koch curve and the Cantor set by repeating a rule, and work
  out what each round does to their length}
\outcome{work out a dimension from how many copies a shape is made of and how
  much each copy is shrunk}
\outcome{compute a box-counting dimension from a list of points, read it off a
  log-log plot, and say how far to trust it}
\outcome{say what a dimension between one and two claims about a shape, and
  over what range of sizes a real object can earn one}
\end{outcomes}

\section{A coastline and a ruler}
\label{sec:ch11-coast}
\index{Richardson, Lewis Fry}\index{Mandelbrot, Benoit}

\noindent
The first person to take the question seriously was Lewis Fry Richardson,
an English physicist. Studying wars late in his life, he wanted to know
whether countries with a longer shared border fought more often, and found
that neighbours could quote quite different lengths for the same border.
His measurements were published in 1961, after his death.

You measure a curve on a map with dividers: two compass points set a fixed
distance apart. Walk them along the coast, count the steps and multiply by
the step. That fixed distance is your ruler.

\begin{think}
You measure a stretch of coast twice on the same detailed map: once with the
dividers set to $10$ kilometres, once set to $1$ kilometre. Which measurement
is longer, or do they agree?
\end{think}
\reveal{The $1$-kilometre measurement. The long ruler steps straight across
every bay narrower than itself; the short one goes into the bay and out
again, and each detour adds length.}
In 1967 Benoit Mandelbrot built a paper in the journal \emph{Science} on
Richardson's measurements, with the title \enquote{How long is the coast of
Britain?}. His answer, which this chapter works towards, was that the
question has none, and that a better question does.

First, see what the experiment does to a curve with no hidden detail: half a
circle of radius one, whose length is $\pi$, about \val{ch11.pi} (half of
the $2\pi$ round a whole circle). Put points on it, walk from each to the
next in a straight line, and add up the steps. That is the whole ruler:

\pyregion{code/ch11/_ruler.py}{length}{The ruler: add up the straight steps}{lst:ch11-length}

\noindent
\code{pairwise(points)} hands the loop each point with the next one, and the
line marked \code{Pythagoras} is the distance between them: across squared
plus up squared, square-rooted. The next listing starts with three points on
the half circle and, round after round, adds a point on the circle halfway
between every pair of neighbours, which halves the ruler:

\pyregion{code/ch11/semicircle.py}{table}{Half a circle, measured with shorter and shorter rulers}{lst:ch11-semicircle}

\transcript{ch11-semicircle}

\noindent
Every shorter ruler measures more, as on the coast. But each round adds about
\val{ch11.semi.gain} of what the round before added, so the additions die
away and the measurements settle on $\pi$. That settling is what a length
\emph{is}: the number your measurements approach as the ruler shrinks. A
smooth curve has one. Richardson's coasts did not: for every coast and
border he tried, the measured length kept climbing as the ruler shrank.

\section{A curve that never settles}
\label{sec:ch11-koch}
\index{Koch curve}\index{self-similarity}

\noindent
You cannot reason about detail you have not yet mapped, so make a coast by
a rule, where every scale of detail is known in advance. Helge von Koch
described this one in 1904. Start with a straight line of length one. Cut it
into three equal pieces and replace the middle one with the two upper sides
of a triangle, each as long as the piece it replaces, like a tent pitched on
the middle third. That is round one. Round two does the same to each of the
four pieces, and so on, for ever.

\begin{think}
Each round replaces every straight piece with four pieces, each a third as
long. By what number does the total length get multiplied each round?
\end{think}
\reveal{By $\tfrac{4}{3}$: four pieces, each $\tfrac{1}{3}$ as long as the
piece they replace, add up to $4 \times \tfrac{1}{3}$ of it.}

\plotfig{ch11-koch-stages}{The Koch curve after rounds $0$ to $5$. Each round
pitches a tent on the middle third of every straight piece; the dashed box,
as tall as the first tent, holds every round.}{the Koch curve, round by round}

\noindent
Look at round $5$: the curve is made of four smaller copies of itself, each
made of four smaller copies again. That is \emph{self-similarity}, the rule
showing through at every scale. The listing builds rounds $0$ to $6$ and
measures each with the ruler of Listing~\ref{lst:ch11-length}:

\pyregion{code/ch11/koch_length.py}{table}{The Koch curve, measured round by round}{lst:ch11-koch}

\transcript{ch11-koch-length}

\noindent
After round $n$ every piece is $\tfrac{1}{3^{n}}$ long, so measuring that
round is measuring the Koch curve with a ruler of that length. The
\code{length} column is $\tfrac{4}{3}$ times the row above, every time: by
round $6$ it is \val{ch11.koch.six}, and by round $20$ it would be
\val{ch11.koch.twenty}. The \code{height} stays at \val{ch11.koch.height},
the first tent's. The \code{area} between the curve and its base (a function
higher in the same file) creeps towards \val{ch11.koch.area} and no
further: each round pitches four times as many tents, each a ninth of the
area, so it adds $\tfrac{4}{9}$ of what the round before added, and those
additions die away like the half circle's.

\begin{projectbox}
The rule is \api{koch} in \pkg{chaoslab}. The one unfamiliar line is the
tent's tip: half a piece along the middle third, then sideways by \code{h}
times the piece, the height of an equilateral triangle.
\pyregion{code/src/chaoslab/fractals.py}{koch}{The Koch rule, in chaoslab}{lst:ch11-kochrule}
\end{projectbox}

\begin{think}
Suppose you could draw every round, for ever, and measure the finished curve
with a ruler as short as you like. What length would you get?
\end{think}
\reveal{No length at all. Each time the ruler shrinks by three the
measurement grows by $\tfrac{4}{3}$, and it never settles. The finished
curve is longer than any number you care to name, and it still fits in the
dashed box.}

\begin{trapbox}
\textbf{\enquote{A coastline has a length; it is only hard to measure.}} A
length is what your measurements settle on as the ruler shrinks, and the
measurements of a coast do not settle: every finer map finds bays inside the
bays. A coastline's length is a fact about the coast \emph{and the ruler}
together, so a figure quoted without its ruler, or the scale of its map, is
not a fact about the coast at all.
\end{trapbox}

\section{Counting copies}
\label{sec:ch11-dimension}
\index{Cantor set}\index{dimension!from copies}

\noindent
Georg Cantor described a second shape in 1883, and it goes the other way.
Take the line from $0$ to $1$, cut it into thirds and throw away the middle
third, keeping its ends. Do the same to each piece that is left, for ever.

\begin{think}
Throw away middle thirds for ever. How much of the line's length is left at
the end? And is anything left at all?
\end{think}
\begin{revealblock}
No length: each round keeps $\tfrac{2}{3}$ of what was there, so after
$20$ rounds \val{ch11.cantor.pieces} pieces hold \val{ch11.cantor.left} of
it, and the total heads for zero. But points are left. The ends of every
piece, such as $\tfrac{1}{3}$ and $\tfrac{2}{3}$, are never thrown away.
Nor is $\tfrac{1}{4}$, which is never an end: it lies in the left third;
stretch that third three times and $\tfrac{1}{4}$ lands on $\tfrac{3}{4}$,
in a right third; stretch that one and it is back at $\tfrac{1}{4}$. It
dodges every middle third for ever.
\end{revealblock}
What is left is the \emph{Cantor set}: a dust of infinitely many points,
more than you could ever write down in a list, holding no length at all.
Length is the wrong question for both shapes: one has too much of it, the
other too little. Here is a better one.

Shrink a line to a third of its size and three copies rebuild the original.
Shrink a square to a third of its width and it takes nine; a cube takes
$27$. Those are $3^{1}$, $3^{2}$ and $3^{3}$, and the power is the
dimension. So a shape made of copies of itself, each shrunk by the same
factor, has the dimension $D$ with
\[ \text{shrink}^{D} = \text{copies}. \]

\begin{think}
The Koch curve is four copies of itself, each shrunk by three. If $3^{D} = 4$,
is $D$ equal to $1$, to $2$, or somewhere between?
\end{think}
\reveal{Between: $3^{1} = 3$ copies is too few and $3^{2} = 9$ is too many.
The power of $3$ that gives $4$ is about \val{ch11.koch.dim}.}
Finding that power is the job of the logarithm, which answered \enquote{how
many steps until} in Chapter~\ref{ch:iteration} and here answers \enquote{how
many factors of three make four}:
\[ D = \frac{\ln 4}{\ln 3} \approx \val{ch11.koch.dim}, \]
where $\ln$ is the natural logarithm on any calculator. The Cantor set is two
copies of itself, each shrunk by three, so its dimension is $\ln 2 / \ln 3$,
about \val{ch11.cantor.dim}: more than separate points, which have dimension
$0$, and less than a line.

\begin{trapbox}
\textbf{\enquote{A dimension must be a whole number.}} It is a whole number
for the shapes you met at school because they are smooth: zoom in on a
smooth curve and it looks like a line. Count copies instead of directions
and dimension becomes a measurement that can come out anywhere. The Koch
curve's \val{ch11.koch.dim} says something precise: too long to have a
length, too thin to have an area. A shape whose dimension is not a whole
number is what Mandelbrot, in 1975, named a \emph{fractal}, from the Latin
\emph{fractus}, broken.
\end{trapbox}

\noindent
Dimension also explains Richardson's rulers. Shrink the ruler by three and
the measured length is multiplied by $3^{D - 1}$: by $\tfrac{4}{3}$ for the
Koch curve, by $1$ for anything smooth, by $\tfrac{2}{3}$ for the Cantor set.
On axes where every step is a factor of ten each is a straight line, and
Figure~\ref{fig:ch11-three-fates} draws all three. Richardson's coasts gave
lines of exactly this kind, and Mandelbrot read dimensions off their slopes:
about \num{1.25} for the west coast of Britain.

\plotfig{ch11-three-fates}{What happens to a measured length as the ruler
shrinks (rightwards). A smooth curve settles; the Koch curve climbs for ever
and the Cantor set shrinks to nothing, each along a straight line whose slope
is its dimension minus one.}{three fates of a length}

\begin{lab}{l11_01_cantor}{The Cantor set by rule}
Write \code{cantor}, which returns the pieces left after a number of rounds,
and \code{length\_left}, which adds them up. Then write
\code{similarity\_dimension}, which turns \enquote{this many copies, each
shrunk by this factor} into a dimension with \code{math.log}. The test checks
the line, the square, the Koch curve and the Cantor set, and that
$\tfrac{1}{4}$ survives every round.
\end{lab}

\section{Counting boxes}
\label{sec:ch11-boxes}
\index{box counting}\index{dimension!box-counting}

\noindent
Counting copies needs a rule, and no rule built a coast. What works on
anything you can draw is \emph{box counting}: lay a grid of squares of side
$s$ over the shape, count the squares it touches, call that $N$, halve $s$
and count again.

\mermaidfig{ch11-box-recipe}{Box counting in three steps. The recipe needs a
list of points and nothing else, so it works on shapes that no rule
built.}{cover, count, slope}

\noindent
For a line, halving the box doubles the count; for a filled square it
multiplies the count by four. In general, halving the box multiplies $N$ by
$2^{D}$, and that is how $D$ is read off.

\begin{think}
If the Koch curve's dimension really is \val{ch11.koch.dim}, about how many
times as many boxes should it need each time the box size is halved?
\end{think}
\reveal{About $2^{\val{ch11.koch.dim}}$, which is \val{ch11.koch.halve}:
more than a line's $2$ and fewer than a square's $4$.}

\pyregion{code/ch11/koch_boxes.py}{count}{Box counting on the Koch curve}{lst:ch11-kochboxes}

\transcript{ch11-koch-boxes}

\noindent
The listing covers the \val{ch11.koch.corners} corners of the round-$8$ curve
with boxes from $\tfrac{1}{2}$ down to $\tfrac{1}{1024}$, in halves on
purpose: the curve is built in thirds, and a measurement that knew that
would be less of a test. The last column wobbles around
\val{ch11.koch.halve}, because thirds and halves do not line up. To average
the wobble away, plot the logarithm of each count against the logarithm of
each $1/s$. If $N$ grows like $(1/s)^{D}$, the logarithm turns that power
into a straight line of slope $D$, and \api{fit\_slope} draws the line that
passes as close as it can to all ten points. Its slope is
\val{ch11.koch.boxdim}; $\ln 4 / \ln 3$ is \val{ch11.koch.dim}.

Do not trust that agreement further than it deserves. Fit through any five
neighbouring sizes instead of all ten and the slope lands anywhere from
\val{ch11.koch.five.lo} to \val{ch11.koch.five.hi}. That spread is the
honest size of the error, and it is why a box-counting dimension is fitted
over as wide a range of sizes as the data allow.

\begin{lab}{l11_02_boxes}{Count the boxes yourself}
Write \code{count\_boxes}, which says how many squares of a given side hold at
least one point, and \code{box\_dimension}, which fits the slope over a list
of sizes. The test tries a straight line, a filled square, the Koch curve and
the Cantor set, whose box-counting dimension should agree with the one you
got from copies in the first lab.
\end{lab}

\section{The shape chaos leaves behind}
\label{sec:ch11-henon}
\index{Hénon attractor!dimension}\index{strange attractor!as fractal}

\noindent
Chapter~\ref{ch:strange} ran Hénon's map, which takes a point $(x, y)$ of the
plane to $(1 - \num{1.4}x^{2} + y,\ \num{0.3}x)$, and found that its points
settle onto an attractor with no area at all, whose stretching gives it no
finite length either. Nobody built that shape copy by copy. It is exactly
the kind of object box counting is for.

\begin{think}
No area and no finite length. Before any box is counted: would you expect the
Hénon attractor's dimension to be $1$, $2$, or something between?
\end{think}
\reveal{Something between $1$ and $2$: no finite length makes it more than a
curve, and no area makes it less than a surface. The listing measures about
\val{ch11.henon.dim}.}

\pyregion{code/ch11/henon_boxes.py}{count}{Box counting on the Hénon attractor}{lst:ch11-henonboxes}

\transcript{ch11-henon-boxes}

\noindent
The function \code{attractor} runs the map from $(0, 0)$ and keeps every
point after the first thousand steps. With
\val{ch11.henon.many} points the slope through all nine box sizes is
\val{ch11.henon.dim}, and through any five neighbouring sizes it lies between
\val{ch11.henon.five.lo} and \val{ch11.henon.five.hi}. With
\val{ch11.henon.few} points it drops to \val{ch11.henon.dim.few}: the
smallest boxes need many points before every thin layer of the attractor
has been visited, so too few points leave boxes uncounted and flatten the
line.

Careful published measurements, with far more points and smaller boxes, put
the dimension at about \num{1.26}. Ours is a little under, has a spread, and
has a known reason to be low: that is what an honest measurement of a
dimension looks like. That the published figure matches the Koch curve's is
a coincidence: the two shapes know nothing of each other.

\plotfig{ch11-box-counts}{Left: the Hénon attractor covered by boxes of side
$\tfrac{1}{16}$. Right: boxes touched against boxes per unit length, on
axes where every step is a factor of ten. Both shapes lie on straight lines
between the dashed lines of slope $1$ and $2$.}{box counts on a log-log plot}

\noindent
This is the sense in which the attractors of this part of the book are
strange. Stretching makes them long, folding keeps them in a box, and a
shape too long for a length and too thin for an area is a fractal.

\section{Fractals in the world}
\label{sec:ch11-world}
\index{fractal!in nature}

\noindent
Once you know the shape you see it everywhere. A fern frond is made of
smaller fronds. A river gathers streams that gather smaller streams. A
romanesco broccoli is a spiral of cones, each a spiral of smaller cones.
Blood vessels branch, and branch again.

\begin{think}
Your lungs must pass oxygen into the blood across a thin surface, and all of
that surface must fit inside your chest. Why might a lung be built like a
fractal, branching and branching again?
\end{think}
\reveal{Because a fractal packs an enormous surface into a small space, as
the Koch curve packs an unlimited length into its dashed box. Every round of
branching multiplies the surface while the chest stays the same size.}

\begin{trapbox}
\textbf{\enquote{Fractals are decoration: pretty pictures, no use.}} Branching
in lungs, blood vessels and river networks does work: it reaches every point
of a volume or a landscape and packs a large surface or a long edge into a
bounded space. The shape is the solution to a problem, and its dimension is a
number you can measure and compare.
\end{trapbox}

\begin{worldbox}
In your lungs the airways divide in two about twenty-three times between the
windpipe and the tiny air sacs where gas crosses into the blood, folding a
surface of several tens of square metres into your chest. Physiologists
describe the airway tree with this chapter's ideas: each generation of
branches is a smaller copy of the one before.
\end{worldbox}

\noindent
Every one of these is fractal only over a range of sizes. The Koch curve has
tents inside tents all the way down, because a rule does not get tired; a
coast stops showing new detail around its pebbles and grains of sand, a fern
at its smallest leaflets, a lung at air sacs a fraction of a millimetre
across. So when somebody calls a real object a fractal, ask: over what range
of sizes, and measured how? A box count spanning a factor of a thousand is a
finding; one spanning a factor of three is a guess.

The most famous fractal of all comes from iterating a rule on the points of
a plane rather than of a line, and it is the subject of
Chapter~\ref{ch:mandelbrot}.

\begin{summarybox}
\begin{itemize}
\item A curve's \textbf{length} is what its measurements settle on as the
  ruler shrinks. A smooth curve settles; a coastline keeps growing, as
  Richardson found and Mandelbrot asked about in 1967.
\item The \textbf{Koch curve} multiplies its length by $\tfrac{4}{3}$ every
  round, so it has no finite length, while its height and the area it bounds
  stay fixed. It is \textbf{self-similar}: made of smaller copies of itself.
\item The \textbf{Cantor set} keeps $\tfrac{2}{3}$ of its length every round,
  so no length is left, yet infinitely many points are. A shape made of
  copies of itself, each shrunk by the same factor, has \textbf{dimension}
  $D$ with $\text{shrink}^{D} = \text{copies}$: about
  \val{ch11.koch.dim} for Koch and \val{ch11.cantor.dim} for Cantor. A
  shape whose dimension is not a whole number is a \textbf{fractal}.
\item \textbf{Box counting} measures a dimension from points alone: the
  slope of $\ln N$ against $\ln(1/s)$, where $N$ boxes of side $s$ touch the
  shape. The spread over different ranges of sizes is the honest error.
\item The \textbf{Hénon attractor} measures about \val{ch11.henon.dim}
  (published: about \num{1.26}): strange attractors are fractals.
\item Real objects \textbf{are fractal over a range of sizes}, never all the
  way down, and fractal branching packs a large surface into a small space.
\end{itemize}
\end{summarybox}

\canyou

\begin{checkpoint}
\item Two atlases give the same coastline as $800$ and $1100$ kilometres long.
  Can both be right, and what should each have told you?
  \answerto{Both can be right if they measured with different rulers (maps of
  different scale): the shorter ruler follows more bays. Each should have
  said what ruler or map scale it used.}
\item A Koch curve starts as a line of length $1$. How long is it after three
  rounds?
  \answerto{$\bigl(\tfrac{4}{3}\bigr)^{3} = \tfrac{64}{27}$, about
  \num{2.37}.}
\item A different rule cuts every piece of a line into five equal parts and
  throws away the second and the fourth. How much length is left after one
  round, and what is the dimension of what is left for ever?
  \answerto{$\tfrac{3}{5}$ of the length. Three copies, each shrunk by five,
  give $D = \ln 3 / \ln 5$, about \val{ch11.five.dim}.}
\item A shape is made of $8$ copies of itself, each shrunk by $2$. What is its
  dimension, and what familiar shape has it?
  \answerto{$2^{3} = 8$, so $D = 3$: a solid cube.}
\item Halving the box size takes a box count from $100$ to $283$. Roughly what
  is the dimension?
  \answerto{About \num{1.5}, because $2^{\num{1.5}} = 2 \times \sqrt{2}$,
  about \num{2.83}.}
\item Why did the Hénon attractor's measured dimension rise when the number of
  points went up?
  \answerto{Too few points leave some of the smallest boxes unvisited, which
  undercounts them and flattens the fitted line. More points fill them in.}
\item Somebody tells you a head of broccoli is a fractal. What should you ask?
  \answerto{Over what range of sizes, and measured how. A real object can be
  fractal only between a largest and a smallest size.}
\end{checkpoint}
